% English translation of questions/5784/semA-special/q3.tex (metadata is taken from the Hebrew file).

\begin{question}
Find the absolute minimum and maximum of the function
\[ f(x)=(x+2)\cdot e^x \]
on its domain, or explain why they do not exist.
\end{question}

\begin{hint}
Find the intervals of increase and decrease using $f'$, and check the behavior of the function as $x\to\pm\infty$.
\end{hint}

\begin{solution}
\textbf{Domain:} all of $\R$; $f$ is differentiable at every point.

\textbf{Derivative:} by the product rule
\[ f'(x)=e^x+(x+2)e^x=(x+3)e^x. \]
Since $e^x>0$, the sign of $f'$ is the sign of $x+3$: $f'<0$ on $(-\infty,-3)$ and $f'>0$ on $(-3,\infty)$. Hence (by a corollary of Lagrange's theorem) $f$ is strictly decreasing on $(-\infty,-3]$ and strictly increasing on $[-3,\infty)$.

\textbf{Absolute minimum:} by monotonicity, for every $x\le-3$ we have $f(x)\ge f(-3)$, and for every $x\ge-3$ we have $f(x)\ge f(-3)$. Hence $x=-3$ is a point of absolute minimum, and the minimum value is
\[ f(-3)=(-3+2)e^{-3}=-\frac1{e^3}\approx-0.0498. \]

\textbf{Absolute maximum:} $\lim_{x\to\infty}(x+2)e^x=+\infty$, hence the function is not bounded above and has no absolute maximum. (For completeness: $\lim_{x\to-\infty}(x+2)e^x=\lim_{x\to-\infty}\frac{x+2}{e^{-x}}=0$ by L'Hôpital's rule, so the function tends to $0^-$ on the left, but on the right it is unbounded.)

\textbf{Answer:} the absolute minimum is $-e^{-3}$, attained at $x=-3$; an absolute maximum does not exist since $f(x)\to+\infty$ as $x\to\infty$.
\end{solution}
